\documentclass[11pt,reqno,a4paper]{amsart}

\usepackage[T1]{fontenc}
\usepackage{lmodern}
\usepackage{amsmath,amssymb,mathtools}
\usepackage{booktabs}
\usepackage{enumitem}
\usepackage{microtype}
\usepackage[hidelinks]{hyperref}

\newtheorem{theorem}{Theorem}[section]
\newtheorem{proposition}[theorem]{Proposition}
\newtheorem{lemma}[theorem]{Lemma}
\newtheorem{corollary}[theorem]{Corollary}
\theoremstyle{definition}
\newtheorem{definition}[theorem]{Definition}
\theoremstyle{remark}
\newtheorem{remark}[theorem]{Remark}

\newcommand{\comp}[1]{\overline{#1}}
\newcommand{\one}{\mathbf{1}}
\newcommand{\Ffive}{\mathbb{F}_{5}}

\title[The five-color case]{The five-color case of an Erd\H{o}s--Gy\'arf\'as balanced-coloring problem}
\author{Robert Sneiderman}
\date{17 July 2026}
\subjclass[2020]{05C15, 05C35, 05D10}
\keywords{edge coloring, balanced coloring, set-coloring Ramsey number, extremal graph theory}

\begin{document}
\hypersetup{
  pdftitle={The five-color case of an Erdos--Gyarfas balanced-coloring problem},
  pdfauthor={Robert Sneiderman}
}

\begin{abstract}
We prove the five-color case of a conjecture of Erd\H{o}s and Gy\'arf\'as: every
five-coloring of the edges of $K_{26}$ has six vertices whose induced edges omit
at least one color.  Equivalently, the set-coloring Ramsey number
$R(6;5,4)$ equals $26$.  The proof assigns a graph to each color, applies the
Kang--Pikhurko extremal theorem to complements of induced subgraphs, and combines
the resulting edge bounds with minimum-degree decompositions.  A sequence of
small near-extremal lemmas first forces all five color graphs to have $65$ edges
and then excludes the equality case.  The argument is non-computational; finite
searches used during its development are not premises of the proof.
\end{abstract}

\maketitle

\begin{center}
\fbox{\parbox{0.88\textwidth}{\centering
\textbf{Preprint.}
This manuscript has not yet completed external mathematical review.}}
\end{center}

\paragraph{AI-use disclosure.}
OpenAI GPT-5.6 Sol was used extensively during proof search, drafting, and
internal checking. OpenAI GPT-5 (Codex) was used for later source review and
release checks. xAI Grok 4.5 was used to critique an earlier draft. This
manuscript has not completed external mathematical review.

\section{Introduction}

Erd\H{o}s and Gy\'arf\'as asked whether the following assertion holds for every
integer $r\geq 3$: whenever the edges of $K_{r^2+1}$ are colored with $r$
colors, some $r+1$ vertices induce a complete graph on which at least one color
is absent.  They proved the cases $r=3$ and $r=4$, observed that the analogous
assertion fails for $r=2$, and studied the related problem on $r^2$ vertices
\cite{ErdosGyarfas1999}.  The problem also appears in Erd\H{o}s's later problem
collection \cite{Erdos1999} and as Erd\H{o}s Problem 617
\cite{ErdosProblems617}.

The purpose of this paper is to prove the assertion for $r=5$.

\begin{theorem}\label{thm:main}
For every map $\chi:E(K_{26})\to[5]$, there is a set
$S\subseteq V(K_{26})$ with $|S|=6$ such that
\[
  \chi\!\left(E(K_{26}[S])\right)\ne[5].
\]
\end{theorem}

This is a fixed-parameter result.  In particular, Theorem~\ref{thm:main} does
not settle the conjecture for all $r\geq 3$.

There is an equivalent formulation in the language of set-coloring Ramsey
numbers.  An $(r,s)$-coloring assigns to each edge an $s$-element subset of a
palette of $r$ colors.  The number $R(n;r,s)$ is the least $N$ for which every
$(r,s)$-coloring of $K_N$ has a copy of $K_n$ whose edges share a common color
\cite{ConlonEtAl2024}.  Replacing the color on each edge by the other four
colors turns Theorem~\ref{thm:main} into the upper bound $R(6;5,4)\leq26$.

The reverse bound has a short affine construction.  Take the vertex set
$\Ffive^2$.  Its six parallel classes of affine lines are indexed by the
slopes $\Ffive\cup\{\infty\}$.  Color a pair of points by the slope of the line
through them, after merging the classes of slopes $0$ and $\infty$ into one
color.  For each of the four unmerged slopes, any six points contain two on a
common line of that slope, by the pigeonhole principle applied to its five
parallel lines.  The same argument with the horizontal lines supplies the
merged color.  Thus every six points see all five colors.  Complementing the
edge colors gives a $(5,4)$-coloring of $K_{25}$ with no monochromatic $K_6$.

\begin{corollary}\label{cor:set-color}
$R(6;5,4)=26$.
\end{corollary}

We now describe the proof of Theorem~\ref{thm:main}.  Assuming a counterexample,
let $G_i$ be the graph formed by color $i$.  Every six-set spans between one
and eleven edges in each $G_i$.  The Kang--Pikhurko theorem first gives a
uniform lower bound on $e(G_i)$ and confines a least color to minimum degree
two, three, or four.  We then prove three layers of induced-subgraph
obstructions, for independence numbers two, three, and four.  The last layer
forces $e(G_i)\geq65$ for every $i$, so all five color graphs have exactly
$65$ edges.  Equality in the minimum-degree counts repeatedly peels off
isolated cliques.  The remaining near-extremal graphs are excluded by two
explicit classifications on fifteen vertices.

\section{Preliminaries}

All graphs are finite and simple.  We write $v(F)=|V(F)|$, $e(F)=|E(F)|$,
$\alpha(F)$ for the independence number, $\omega(F)$ for the clique number,
$\tau(F)$ for the minimum vertex-cover number, and $\delta(F)$ and
$\Delta(F)$ for the minimum and maximum degrees.  The
complement of $F$ is $\comp F$, and $F[W]$ is the subgraph induced by $W$.
The Tur\'an graph with $r$ parts on $n$ vertices is $T_r(n)$, and
$t_r(n)=e(T_r(n))$.

Suppose, for the rest of the proof, that a counterexample to
Theorem~\ref{thm:main} is given.  For a fixed color, let $G$ be its color
graph.  For every $S\in\binom{V(G)}6$,
\begin{equation}\label{eq:local}
  1\leq e(G[S])\leq11.
\end{equation}
Indeed, the lower bound says that the chosen color is present on $S$, while
the upper bound follows because each of the other four colors occupies at
least one of the fifteen edges.  In particular,
\begin{equation}\label{eq:alpha-omega}
  \alpha(G)\leq5,
  \qquad
  \omega(G)\leq5.
\end{equation}
Every induced subgraph of a color graph retains the upper bound in
\eqref{eq:local}.

\begin{definition}
A graph is \emph{admissible} if every six vertices span at most eleven edges.
\end{definition}

The main external extremal result is the following theorem of Kang and
Pikhurko \cite[Theorems 1 and 4, Lemma 5]{KangPikhurko2005}.

\begin{theorem}[Kang--Pikhurko]\label{thm:KP}
Let $n\geq r+3$, $r\geq2$, and let $L$ be a $K_{r+1}$-free graph on $n$
vertices that is not $r$-partite.  If $r\leq(n-1)/2$, then
\begin{equation}\label{eq:KP}
  e(L)\leq t_r(n)-\left\lfloor\frac nr\right\rfloor+1.
\end{equation}
Equality holds only for the graphs in the construction described
below.
\end{theorem}

For the equality construction, choose positive integers
$n_1\leq\cdots\leq n_r$ with sum $n-1$ and at least two entries larger than
one.  Let $N_1,\ldots,N_r$ be disjoint sets of these sizes, and let $s,t$, in
either order, be the two smallest indices whose parts have size greater than
one.  Choose a proper nonempty set $A\subset N_s$ and a vertex $y\in N_t$.
Begin with the
complete $r$-partite graph on the $N_i$, add a vertex $x$ adjacent to
\[
  \bigcup_{i\notin\{s,t\}}N_i\;\cup\;A\;\cup\;\{y\},
\]
and delete all edges between $y$ and $A$.  Kang and Pikhurko prove that all
equality graphs arise this way, with a short list of optimal part sizes.
The instances needed below are
\begin{center}
\begin{tabular}{ccl}
\toprule
$r$ & $n$ & optimal sequences $(n_1,\ldots,n_r)$\\
\midrule
2 & 10 & $(4,5)$\\
3 & 15 & $(4,4,6)$, $(4,5,5)$\\
\bottomrule
\end{tabular}
\end{center}

We will repeatedly exclude equality in \eqref{eq:KP} using
\eqref{eq:local}.  The required checks are recorded once here.

\begin{lemma}[Local exclusion of equality]\label{lem:KP-equality}
Let $R=\comp L$ be admissible.
\begin{enumerate}[label=\textup{(\roman*)}]
\item If $(r,n)=(2,10)$ and $L$ is nonbipartite, then $e(L)\leq20$.
\item If $(r,n)=(3,15)$ and $L$ is not three-partite, equality in
\eqref{eq:KP} is impossible.
\end{enumerate}
\end{lemma}

\begin{proof}
For (i), the unique sequence is $(4,5)$.  If the part of size five is the
$t$-part, that part together with $x$ spans fourteen edges in $R$.  If it is
the $s$-part and $|A|\leq3$, the same six-set spans at least twelve edges.  If
$|A|=4$, use the size-five part together with $y$; it spans fourteen edges.

For (ii), the sequence $(4,4,6)$ gives a $K_6$ in $R$.  Consider
$(4,5,5)$.  If a size-five $t$-part is used, that part with $x$ spans
fourteen edges in $R$.  If a size-five $s$-part is used and $|A|\leq3$,
that part with $x$ spans at least twelve edges.  If $|A|=4$, the size-five
$s$-part with $y$ spans fourteen edges.
\end{proof}

We also use Brooks's theorem in the standard form: a connected graph of
maximum degree $D$ is $D$-colorable unless it is complete or an odd cycle
\cite{Brooks1941}.

\begin{lemma}[Neighborhood accounting]\label{lem:accounting}
Let $F$ be a graph, let $v$ be a minimum-degree vertex of degree $d$, let
$N=N_F(v)$, and let $U=V(F)\setminus(N\cup\{v\})$.  Put
\[
  X=e_F(N,U),\qquad Y=e(F[N]),\qquad A=X+Y.
\]
Then
\begin{equation}\label{eq:accounting}
  X+2Y\geq d(d-1),\qquad
  Y\leq\binom d2,\qquad
  A\geq\binom d2.
\end{equation}
If $F$ is admissible and $d=5$, then $Y\leq6$ and $A\geq14$.  If equality
$A=\binom d2$ holds in the last inequality of \eqref{eq:accounting}, then
$N\cup\{v\}$ is an isolated $K_{d+1}$ in $F$.
\end{lemma}

\begin{proof}
Every vertex of $N$ has at least $d-1$ incident edges other than its edge to
$v$.  Summing these requirements over $N$ counts the $N$--$U$ edges once and
the edges inside $N$ twice, proving the first inequality.  Combining it with
$Y\leq\binom d2$ gives the third.  When $d=5$, the six-set
$N\cup\{v\}$ has $5+Y\leq11$ edges.  Finally,
$X+Y=\binom d2$ and $X+2Y\geq d(d-1)=2\binom d2$ force
$Y=\binom d2$ and $X=0$.
\end{proof}

\begin{proposition}[Initial structural reduction]\label{prop:initial}
If $G$ is a least frequent color graph, then
\[
  e(G)\leq65,
  \qquad
  2\leq\delta(G)\leq4.
\]
\end{proposition}

\begin{proof}
The five edge counts sum to $\binom{26}{2}=325$, so a least color has at
most $65$ edges.  If its minimum degree were at least five, its average
degree would force it to be five-regular.  Each component is neither $K_6$
nor an odd cycle: $K_6$ is excluded by \eqref{eq:alpha-omega}, and an odd
cycle is not five-regular.  Brooks's theorem gives a proper five-coloring.  One
color class would have at least six vertices, contrary to
$\alpha(G)\leq5$.  Hence $\delta(G)\leq4$.

Let $v$ be a minimum-degree vertex of degree $d$.  Its $25-d$ nonneighbors
induce a graph $G[U]$ with independence number at most four.  Thus
$\comp{G[U]}$ is $K_5$-free.  When $d=0$ or $1$, this complement is not
four-partite: a four-partition of $25$ or $24$ vertices has a part of size
at least six, which would be a forbidden clique in $G[U]$.  For $d=0$,
Theorem~\ref{thm:KP} gives $e(G[U])\geq71$.  For $d=1$, the same theorem on
24 vertices gives $e(G[U])\geq65$, and the edge incident with $v$ gives
$e(G)\geq66$.  Both cases contradict $e(G)\leq65$, so $d\geq2$.
\end{proof}

\section{Terminal bounds at independence number two}\label{sec:terminal}

The first layer concerns complements of triangle-free graphs.  Notice that
if an admissible graph $F$ has $\alpha(F)\leq2$, then $L=\comp F$ is
triangle-free and every six vertices span at least four edges in $L$.

\begin{lemma}\label{lem:alpha2-large}
Let $F$ be admissible with $\alpha(F)\leq2$ and $v(F)=n\geq12$.  Then
\[
  e(F)\geq\binom n2-2n.
\]
\end{lemma}

\begin{proof}
Put $L=\comp F$.  A vertex of degree at least six in $L$ has six independent
neighbors, a contradiction.  Suppose that $w$ has degree five.  Let
$A=N_L(w)$ and let $B$ be the remaining $n-6$ vertices.  The set $A$ is
independent.  For each $x\in B$, the six-set $A\cup\{x\}$ forces
$d_A(x)\geq4$.  On the other hand, each vertex of $A$ uses one incident edge
on $w$ and has total degree at most five, so
\[
  e_L(A,B)\leq5\cdot4=20.
\]
Since $n\geq12$, the lower bound $4|B|\geq24$ is impossible.  Therefore
$\Delta(L)\leq4$ and $e(L)\leq2n$.
\end{proof}

The order eleven endpoint needs a sharper argument.

\begin{lemma}\label{lem:alpha2-eleven}
Let $F$ be an admissible graph on eleven vertices with $\alpha(F)\leq2$.
Then $e(F)\geq36$.
\end{lemma}

\begin{proof}
Put $L=\comp F$.  As above, $L$ is triangle-free, every six-set spans at
least four $L$-edges, and $\Delta(L)\leq5$.  A degree-five vertex is also
impossible.  Indeed, if $A=N_L(w)$ and $B$ is the other five vertices, each
$x\in B$ has at least four neighbors in $A$.  Two vertices of $B$ cannot be
adjacent, because their two subsets of at least four neighbors in the
five-set $A$ intersect and would form a triangle.  Thus $B\cup\{w\}$ is an
independent six-set in $L$, contrary to the four-edge condition.  Hence
$\Delta(L)\leq4$.

Suppose that $e(F)\leq35$, so $e(L)\geq20$.  Choose a degree-four vertex
$w$, put $A=N_L(w)$, and let $B$ be the remaining six vertices.  For
$b\in B$, write $p_b=d_A(b)$.  For distinct $x,y\in B$, the six-set
$A\cup\{x,y\}$ gives
\begin{equation}\label{eq:pair-p}
  p_x+p_y+\one_{xy\in E(L)}\geq4.
\end{equation}
Also,
\begin{equation}\label{eq:sum-p}
  \sum_{b\in B}p_b=e_L(A,B)\leq4\cdot3=12.
\end{equation}
These inequalities force $p_b=2$ for every $b\in B$.  To verify this, order
the six values.  Two values at most one violate \eqref{eq:pair-p}; a zero
forces the other five values to be at least three, contradicting
\eqref{eq:sum-p}.  If one value is one, all others are at least two.  The
only remaining totals at most twelve are $(1,2,2,2,2,2)$ and
$(1,2,2,2,2,3)$.  The vertex with value one must then be adjacent in $B$ to
at least four vertices with value two, although its $B$-degree is at most
three.

It follows that every vertex of $A$ has degree four and every vertex of $B$
has $B$-degree at most two.  Also
\[
  e(L[B])=e(L)-4-12.
\]
If $e(L)\geq23$, this exceeds the maximum possible number of edges in
$L[B]$.  For $e(L)=20,21,22$, the graph $L[B]$ has respectively four,
five, or six edges.  Along an edge of $L[B]$, the two 2-subsets of $A$ given
by the endpoint neighborhoods are disjoint, hence complementary.

If $L[B]$ has six edges, triangle-freeness and maximum degree two force
$L[B]=C_6$.  Its two parity classes carry alternating complementary labels
$P$ and $A\setminus P$.  The set consisting of $w$, the two vertices of
$P$, and the three vertices in the parity class labeled $A\setminus P$
spans only the two edges from $w$ to $P$.

If $L[B]$ has five edges, its triangle-free possibilities are $P_6$,
$C_5$ plus an isolated vertex, and $C_4$ plus $P_2$.  The odd cycle is
incompatible with alternating complementary labels.  The path gives the
same parity-class contradiction as $C_6$.  In the last case, take two
independent cycle vertices with one label.  Of the endpoints of $P_2$,
choose one whose label meets the complementary 2-set in at most one point.
Together with $w$ and that complementary 2-set, these vertices span at most
three edges.

If $L[B]$ has four edges, its triangle-free possibilities are $P_5$ plus an
isolated vertex, $P_4$ plus $P_2$, two copies of $P_3$, and $C_4$ plus two
isolated vertices.  The path case is again the parity argument.  For the
next two cases, take two independent same-label vertices in one path and a
vertex from the other component whose label meets the complementary 2-set
in at most one point.  In the $C_4$ case, use two opposite same-label cycle
vertices and a suitable isolate.  If neither isolate is suitable, both
carry the complementary label, and the two isolates with a cycle vertex of
that label give the required set.  In every case there is a six-set with at
most three $L$-edges, a contradiction.
\end{proof}

We will also need the structure at order ten.

\begin{lemma}[Ten-vertex structure]\label{lem:ten-structure}
Let $L$ be triangle-free on ten vertices, and suppose every six vertices
span at least four edges.
\begin{enumerate}[label=\textup{(\roman*)}]
\item If $L$ is bipartite, its parts both have size five.  In $\comp L$ the
two parts are $K_5$'s and the cross-edges form a matching.
\item If $L$ is nonbipartite, then $e(L)\leq20$.
\item If $L$ is nonbipartite and $e(L)=20$, then $L$ is the balanced
two-fold blow-up of $C_5$.
\item If $e(L)=19$, then $\alpha(L)\leq4$.
\end{enumerate}
For the balanced two-fold blow-up $T$, one has $\alpha(T)=4$ and
$\tau(T)=6$.  Its minimum vertex covers are unions of three of the five
parts, and two distinct minimum covers intersect in at most four vertices.
\end{lemma}

\begin{proof}
If $L$ is bipartite, each part has size at most five, since an independent
six-set is forbidden.  Thus both parts have size five.  A vertex of one
$K_5$ in $\comp L$ has at most one neighbor in the other, since that vertex
together with the first $K_5$ would otherwise span at least twelve
complement edges.  This proves (i).

For nonbipartite $L$, Theorem~\ref{thm:KP} gives $e(L)\leq21$, and
Lemma~\ref{lem:KP-equality}(i) excludes equality.  This proves (ii).

Suppose $e(L)=20$.  A degree-five vertex would have an independent
five-neighborhood $A$.  Each of the other four vertices would need at least
four neighbors in $A$, and triangle-freeness would make those four vertices
independent.  This would require at least $5+16=21$ edges.  Hence $L$ is
four-regular.  Take a shortest odd cycle $C$ of length $\ell$.  It is
induced, and the $2\ell$ edges from $C$ to its complement fit into only
$4(10-\ell)$ outside degree slots, so $\ell\leq6$.  Triangle-freeness and
oddness give $\ell\geq5$ odd, and hence $\ell=5$.  Each outside vertex has
exactly two nonconsecutive neighbors on $C$.  If $x_i$ counts the outside
vertices adjacent to cycle vertices $i-1$ and $i+1$, the external degree two
of cycle vertex $j$ gives
\[
  x_{j-1}+x_{j+1}=2.
\]
The recurrence around the odd cycle forces $x_i=1$ for every $i$.  The
outside graph is two-regular and triangle-free, hence another $C_5$, and
its allowed edges join consecutive types.  Pair cycle vertex $i$ with the
unique outside vertex of type $i$.  Each pair is independent, and every two
consecutive pairs are joined completely.  Four-regularity leaves no other
edges, so these five pairs exhibit $L$ as the balanced two-fold blow-up of
$C_5$.  This proves (iii).

For (iv), an independent five-set would force each of the remaining five
vertices to have at least four neighbors in it, already accounting for
twenty edges.  Finally, in the balanced blow-up the maximum independent
sets are precisely the unions of two nonconsecutive size-two parts.  The
claims about covers follow by taking complements.
\end{proof}

\section{Obstructions at independence number three}\label{sec:alpha3}

The terminal estimates propagate to the next independence level.

\begin{lemma}\label{lem:alpha3-obstruction}
There is no admissible graph $F$ with $\alpha(F)\leq3$ in any of the
following ranges:
\[
  (v(F),e(F))\in
  \{(16,{\leq}44),(17,{\leq}49),(18,{\leq}53)\}.
\]
\end{lemma}

\begin{proof}
First suppose that $v(F)=16$ and $e(F)\leq44$.  A minimum-degree vertex
$v$ has degree $d\leq5$.  For $d\leq4$, the graph on its $15-d$
nonneighbors has independence number at most two.  Lemmas
\ref{lem:alpha2-large} and \ref{lem:alpha2-eleven}, together with
Lemma~\ref{lem:accounting}, give the lower bounds
\begin{center}
\begin{tabular}{c@{\qquad}ccccc}
\toprule
$d$ & 0 & 1 & 2 & 3 & 4\\
\midrule
$e(F)$ & $75$ & $63+1$ & $52+2+1$ & $42+3+3$ & $36+4+6$\\
\bottomrule
\end{tabular}
\end{center}
and every entry exceeds $44$.

It remains to consider $d=5$.  Use the notation of
Lemma~\ref{lem:accounting}, put $R=F[U]$, and put $L=\comp R$.  Then
\begin{equation}\label{eq:alpha3-ten}
  e(R)\leq44-5-14=25,
  \qquad e(L)\geq20.
\end{equation}
If $L$ is bipartite, Lemma~\ref{lem:ten-structure}(i) says that $R$ consists
of two $K_5$'s with at most five cross-edges, which form a matching.  A
vertex of $N$ has at most one $F$-neighbor in either clique.  If two
vertices of $N$ were nonadjacent, each clique would contain three vertices
avoiding both, and some pair across the two three-sets would be a nonedge of
$R$.  This would give an independent four-set in $F$.  Thus $N=K_5$, and
$N\cup\{v\}$ is a forbidden $K_6$.

Suppose instead that $L$ is nonbipartite.  Lemma
\ref{lem:ten-structure}(ii) and \eqref{eq:alpha3-ten} force equality
throughout:
\[
  e(F)=44;\quad e(R)=25;\quad e(L)=20;\quad
  A=14;\quad Y=6;\quad X=8.
\]
By Lemma~\ref{lem:ten-structure}(iii), $L$ is the balanced two-fold blow-up
$T$ of $C_5$.  Let $M=\comp{F[N]}$.  Then $e(M)=4$.  If $x_a$ denotes the
number of $F$-neighbors of $a\in N$ in $U$, minimum degree gives
$x_a\geq d_M(a)$.  Since both sides sum to eight, equality holds for every
$a$.

For an edge $aa'\in E(M)$, the two $U$-neighborhoods must cover every edge
of $T$; otherwise an uncovered edge of $T$, together with $a,a'$, is an
independent four-set in $F$.  On the other hand,
\[
  x_a+x_{a'}=d_M(a)+d_M(a')\leq e(M)+1=5,
\]
whereas $\tau(T)=6$ by Lemma~\ref{lem:ten-structure}.  This contradiction
settles the sixteen-vertex case.

For the remaining orders the same minimum-degree decomposition is purely
numerical.  The complete tables are
\begin{center}\small
\begin{tabular}{c@{\quad}rrrrrr}
\toprule
& $d=0$ & $d=1$ & $d=2$ & $d=3$ & $d=4$ & $d=5$\\
\midrule
$v(F)=17$ & $88$ & $75+1$ & $63+2+1$ & $52+3+3$
            & $42+4+6$ & $36+5+14$\\
$v(F)=18$ & $102$ & $88+1$ & $75+2+1$ & $63+3+3$
            & $52+4+6$ & $42+5+14$\\
\bottomrule
\end{tabular}
\end{center}
Every entry in the first row exceeds $49$, and every entry in the second
exceeds $53$.
\end{proof}

We also record four direct consequences of Theorem~\ref{thm:KP}.

\begin{lemma}\label{lem:alpha3-large}
If $F$ is admissible, $\alpha(F)\leq3$, and $19\leq v(F)\leq22$, then
\begin{center}
\begin{tabular}{c@{\qquad}rrrr}
\toprule
$v(F)$ & 19 & 20 & 21 & 22\\
\midrule
$e(F)$ & $\geq56$ & $\geq62$ & $\geq69$ & $\geq76$\\
\bottomrule
\end{tabular}
\end{center}
\end{lemma}

\begin{proof}
Put $L=\comp F$.  It is $K_4$-free.  It is not three-partite, since a part
would have at least six vertices and would give a $K_6$ in $F$.  Subtracting
the bound \eqref{eq:KP} from $\binom{v(F)}2$ gives the four displayed
values.
\end{proof}

\section{Near-extremal graphs on fifteen vertices}\label{sec:fifteen}

Two small structure statements will be used in the last step.  In both,
the assumption that $\comp R$ is not three-partite is equivalent to saying
that $V(R)$ cannot be covered by three cliques.

\begin{lemma}\label{lem:fifteen-35}
Let $R$ be an admissible graph on fifteen vertices with $\alpha(R)\leq3$,
and suppose that $\comp R$ is not three-partite.  If $e(R)=35$, then
\[
  R=K_5\mathbin{\dot\cup}R_0,
\]
where $R_0$ has ten vertices and $\comp{R_0}$ is the balanced two-fold
blow-up of $C_5$.
\end{lemma}

\begin{proof}
The average degree is less than five.  A minimum-degree vertex cannot have
degree at most three: the graph on its nonneighbors has independence number
at most two, and for degrees zero, one, two, and three the bounds from
Section~\ref{sec:terminal}, including the vertex and neighborhood
contributions, are respectively
\[
  63,\qquad 52+1,\qquad 42+2+1,\qquad 36+3+3,
\]
all larger than $35$.  Hence $\delta(R)=4$.

Choose a degree-four vertex $w$, let $N=N_R(w)$, and let $Q$ be its ten
nonneighbors.  Put $X=e_R(N,Q)$, $Y=e(R[N])$, and $B=X+Y$.  Then
\begin{equation}\label{eq:R35-split}
  e(R[Q])+B=31,\qquad B\geq6,\qquad e(R[Q])\leq25.
\end{equation}
Also $\alpha(R[Q])\leq2$, since an independent triple in $Q$ together with
$w$ would be an independent four-set.  Thus $T=\comp{R[Q]}$ is
triangle-free and every six vertices span at least four $T$-edges.

Suppose $e(R[Q])\leq24$, so $e(T)\geq21$.  If $T$ is nonbipartite,
Theorem~\ref{thm:KP} and Lemma~\ref{lem:KP-equality}(i) give $e(T)\leq20$,
a contradiction.  Hence $T$ is bipartite.  Its parts both have size five,
and $R[Q]$ consists of two $K_5$'s with at most four matching cross-edges.
Each vertex of $N$ has at most one neighbor in either clique.  If two
vertices of $N$ were nonadjacent, three vertices in each clique would avoid
both, and some cross-pair among them would be a nonedge of $R$.  This gives
an independent four-set.  Therefore $N=K_4$, and $R$ is covered by the
three cliques $N\cup\{w\}$ and the two cliques in $Q$, contrary to the
hypothesis.

It follows from \eqref{eq:R35-split} that $e(R[Q])=25$ and $B=6$.
Lemma~\ref{lem:accounting} now gives $Y=6$ and $X=0$, so $N\cup\{w\}$ is
an isolated $K_5$.  The graph $T$ is not bipartite, since otherwise the two
cliques in $Q$ and this $K_5$ would cover $R$.  It has twenty edges, and
Lemma~\ref{lem:ten-structure}(iii) identifies it as the balanced two-fold
blow-up of $C_5$.
\end{proof}

\begin{lemma}\label{lem:fifteen-36}
Let $R$ be an admissible graph on fifteen vertices with $\alpha(R)\leq3$,
and suppose that $\comp R$ is not three-partite.  If $e(R)=36$, then exactly
one of the following holds.
\begin{enumerate}[label=\textup{(\roman*)}]
\item The graph $R$ contains a $K_5$ joined by exactly one edge to a
ten-vertex graph $R_0$, and $\comp{R_0}$ is the balanced two-fold blow-up
of $C_5$.
\item The graph $R$ is an isolated $K_5$ together with a ten-vertex graph
$R_0$ such that $\comp{R_0}$ is triangle-free, has nineteen edges, and has
independence number at most four.
\end{enumerate}
\end{lemma}

\begin{proof}
As in the proof of Lemma~\ref{lem:fifteen-35}, the minimum degree is four.
With the same notation,
\begin{equation}\label{eq:R36-split}
  e(R[Q])+B=32,\qquad B\geq6,\qquad e(R[Q])\leq26.
\end{equation}
The values $e(R[Q])\leq24$ give the same three-clique-cover contradiction.

Suppose first that $e(R[Q])=25$.  The triangle-free graph
$T=\comp{R[Q]}$ has twenty edges.  If it were bipartite, the same candidate
argument would force $N=K_4$ and again give a three-clique cover.  Hence
$T$ is the balanced two-fold blow-up of $C_5$.  Here $B=7$, and
\[
  B+Y=X+2Y\geq12,
\]
so $Y\geq5$.  If $Y=5$, the graph $R[N]$ is $K_4$ minus one edge.  The two
nonadjacent endpoints have at most $X=2$ neighbors in $Q$, but their
$Q$-neighborhoods must cover every edge of $T$ to avoid an independent
four-set.  This contradicts $\tau(T)=6$.  Therefore $Y=6$ and $X=1$,
which is alternative (i).

Finally suppose that $e(R[Q])=26$.  Then $B=6$, and
Lemma~\ref{lem:accounting} gives $Y=6$ and $X=0$.  Thus $N\cup\{w\}$ is an
isolated $K_5$.  The graph $T$ has nineteen edges.  It cannot be bipartite:
two parts of size five would make $R[Q]$ two $K_5$'s with six cross-edges,
although those cross-edges form a matching and hence number at most five.
Lemma~\ref{lem:ten-structure}(iv) gives $\alpha(T)\leq4$.  This is
alternative (ii).
\end{proof}

\section{Obstructions at independence number four}\label{sec:alpha4}

\begin{lemma}\label{lem:alpha4-obstruction}
There is no admissible graph $F$ with $\alpha(F)\leq4$ in any of the
following ranges:
\[
  (v(F),e(F))\in
  \{(23,{\leq}61),(22,{\leq}58),(21,{\leq}54)\}.
\]
\end{lemma}

\begin{proof}
In each range the average degree is less than six, so a minimum-degree
vertex has degree at most five.  Its nonneighbors induce a graph with
independence number at most three.  Lemmas \ref{lem:alpha3-obstruction} and
\ref{lem:alpha3-large}, together with Lemma~\ref{lem:accounting}, give the
following complete lower-bound table for the first two orders:
\begin{center}\small
\begin{tabular}{c@{\quad}rrrrrr}
\toprule
& $d=0$ & $d=1$ & $d=2$ & $d=3$ & $d=4$ & $d=5$\\
\midrule
$v(F)=23$ & $76$ & $69+1$ & $62+2+1$ & $56+3+3$
            & $54+4+6$ & $50+5+14$\\
$v(F)=22$ & $69$ & $62+1$ & $56+2+1$ & $54+3+3$
            & $50+4+6$ & $45+5+14$\\
\bottomrule
\end{tabular}
\end{center}
Every entry in the first row exceeds $61$, and every entry in the second
exceeds $58$.

For $v(F)=21$, degrees zero through four are also excluded numerically.  The
closest three bounds are
\[
  54+2+1=57,\qquad 50+3+3=56,\qquad 45+4+6=55
\]
for degrees two, three, and four respectively.  It remains to consider
$d=5$.  Let $v$ be such a vertex, let $N=N_F(v)$, let $U$ be its fifteen
nonneighbors, put $R=F[U]$, and put $L=\comp R$.  Then
\begin{equation}\label{eq:alpha4-fifteen}
  e(R)\leq54-5-14=35.
\end{equation}

Suppose first that $L$ is three-partite.  Each part has size five: a part
is a clique in $R$, and admissibility forbids a clique of size six.  Thus
$R$ consists of three $K_5$'s and at most five cross-edges.  Every vertex
of $N$ has at most one neighbor in each clique.  If $a,a'\in N$ were
nonadjacent, each clique would supply at least three vertices avoiding both.
There are $3^3=27$ triples choosing one such vertex from each clique, while
each cross-edge rules out at most three triples.  Since $5\cdot3<27$, an
independent triple in $R$ remains.  Together with $a,a'$ it is an
independent five-set in $F$.  Therefore $N=K_5$, and $N\cup\{v\}$ is a
forbidden $K_6$.

Hence $L$ is not three-partite.  By Theorem~\ref{thm:KP}, we have
$e(R)\geq34$.  Lemma~\ref{lem:KP-equality}(ii) excludes equality.  Thus
$e(R)=35$, and equality holds in all remaining counts:
\[
  e(F)=54;\quad A=14;\quad Y=6;\quad X=8.
\]
Let $M=\comp{F[N]}$.  Then $e(M)=4$, and, as in the proof of
Lemma~\ref{lem:alpha3-obstruction}, each $a\in N$ has exactly $d_M(a)$
neighbors in $U$.

By Lemma~\ref{lem:fifteen-35},
\[
  R=C\mathbin{\dot\cup}R_0,
\]
where $C=K_5$ and $T=\comp{R_0}$ is the balanced two-fold blow-up of $C_5$.
Choose an edge $aa'\in E(M)$.  Each endpoint has at most one neighbor in
$C$, so some $c\in C$ avoids both.  If an edge of $T$ avoids the
$R_0$-neighborhoods of $a,a'$, that edge together with $a,a',c$ is an
independent five-set in $F$.  Thus those two neighborhoods cover $T$, but
their total size is at most
\[
  d_M(a)+d_M(a')\leq e(M)+1=5<\tau(T),
\]
a contradiction.
\end{proof}

\section{Edge equalization and the final reduction}\label{sec:equalization}

The preceding obstruction removes the need to analyze separate cases with
sixty, sixty-one, sixty-two, sixty-three, and sixty-four edges.

\begin{proposition}\label{prop:equalization}
Every color graph has exactly $65$ edges.
\end{proposition}

\begin{proof}
Let $G$ be a least frequent color graph.  By
Proposition~\ref{prop:initial}, $e(G)\leq65$ and a minimum-degree vertex has
degree $d\in\{2,3,4\}$.  With the notation of
Lemma~\ref{lem:accounting},
\begin{equation}\label{eq:nonneighbor-upper}
  e(G[U])=e(G)-d-A
  \leq e(G)-d-\binom d2.
\end{equation}
If $e(G)\leq64$, the right side is at most $61$, $58$, or $54$ for
$d=2$, $3$, or $4$.  The graph $G[U]$ has respectively $23$, $22$, or
$21$ vertices, is admissible, and has independence number at most four.
Lemma~\ref{lem:alpha4-obstruction} excludes all three possibilities.  Hence
$e(G)=65$.  Since the five edge counts sum to
$\binom{26}{2}=325=5\cdot65$, every color graph has $65$ edges.
\end{proof}

Every color is now tied for least frequent, so Proposition
\ref{prop:initial} applies to each color graph.

\begin{lemma}\label{lem:no-degree-two-three}
No color graph has minimum degree two or three.
\end{lemma}

\begin{proof}
Let $G$ be a color graph.  Suppose first that $G$ has a degree-two vertex.
Equation~\eqref{eq:nonneighbor-upper} and
Lemma~\ref{lem:alpha4-obstruction} force
\[
  e(G[U])=62,\qquad A=1,\qquad Y=1,\qquad X=0.
\]
Thus $G=K_3\mathbin{\dot\cup}H$, where $H$ has $23$ vertices, $62$ edges,
minimum degree at least two, and independence number at most four.

The average degree of $H$ is less than six.  At a minimum-degree vertex of
degree two, four, or five, the graph on its nonneighbors has independence
number at most three, and Lemmas \ref{lem:alpha3-obstruction} and
\ref{lem:alpha3-large} give respective total lower bounds
\[
  62+2+1,\qquad 54+4+6,\qquad 50+5+14,
\]
all larger than $62$.  Therefore $\delta(H)=3$.  The nineteen nonneighbors
then induce at least $56$ edges, so equality is forced:
\[
  e(H[U])=56,\qquad A=3,\qquad Y=3,\qquad X=0.
\]
Peeling off the resulting isolated $K_4$ leaves a graph $Q$ on nineteen
vertices with $56$ edges, minimum degree at least three, and independence
number at most three.  Its average degree is less than six.  At a
minimum-degree vertex of degree three, four, or five, the nonneighbor graph
has independence number at most two.  Lemma~\ref{lem:alpha2-large} and the
neighborhood contributions give the respective totals
\[
  75+3+3,\qquad 63+4+6,\qquad 52+5+14.
\]
Every total exceeds $e(Q)=56$.

Now suppose that $\delta(G)=3$.  The same argument gives
\[
  e(G[U])=59,\qquad A=3,\qquad Y=3,\qquad X=0,
\]
so $G=K_4\mathbin{\dot\cup}H$, where $H$ has $22$ vertices, $59$ edges,
and minimum degree at least three.  Its minimum degree is at most five.  If
it is three, the eighteen nonneighbors induce at most
$59-3-3=53$ edges; if it is four, the seventeen nonneighbors induce at
most $59-4-6=49$ edges.  Both possibilities contradict
Lemma~\ref{lem:alpha3-obstruction}.  If it is five, the sixteen
nonneighbors induce at most $59-5-14=40$ edges, another contradiction to
that lemma.
\end{proof}

It follows that every color graph has minimum degree four.  Choosing a
degree-four vertex and applying Lemma~\ref{lem:alpha4-obstruction} once
more forces equality in \eqref{eq:nonneighbor-upper}:
\begin{equation}\label{eq:K5H}
  G=K_5\mathbin{\dot\cup}H,
\end{equation}
where $H$ has twenty-one vertices, $55$ edges, minimum degree at least four,
independence number at most four, and is admissible.  Its average degree is
less than six, so $\delta(H)$ is four or five.

\subsection*{Excluding minimum degree four in the residual graph}

Suppose that $\delta(H)=4$.  At a degree-four vertex,
Lemma~\ref{lem:alpha3-obstruction} forces equality in the neighborhood
decomposition.  Thus
\[
  H=K_5\mathbin{\dot\cup}Q,
\]
where $Q$ has sixteen vertices, $45$ edges, minimum degree at least four,
and independence number at most three.  If $\delta(Q)=4$, the eleven
nonneighbors of a degree-four vertex induce at most $45-4-6=35$ edges and
have independence number at most two, contrary to
Lemma~\ref{lem:alpha2-eleven}.  Hence $\delta(Q)=5$.

Choose a degree-five vertex $v$ of $Q$, let $N=N_Q(v)$, let $U$ be its ten
nonneighbors, put $R=Q[U]$, and put $L=\comp R$.  Then $L$ is triangle-free
and
\begin{equation}\label{eq:Q-ten}
  e(R)\leq45-5-14=26,
  \qquad e(L)\geq19.
\end{equation}
If $L$ is bipartite, Lemma~\ref{lem:ten-structure}(i) says that $R$ consists
of two $K_5$'s with matching cross-edges.  Each vertex of $N$ has at most
one neighbor in either clique.  Two nonadjacent vertices of $N$ would then
give an independent four-set in $Q$, as in the proof of
Lemma~\ref{lem:alpha3-obstruction}.  Thus $N=K_5$, making a forbidden
$K_6$ with $v$.  Therefore $L$ is nonbipartite, and
Lemma~\ref{lem:ten-structure}(ii) leaves only $e(L)=19$ or $20$.

If $e(L)=19$, then $e(R)=26$, and all counts are equalities:
\[
  A=14,\qquad Y=6,\qquad X=8.
\]
Let $M=\comp{Q[N]}$.  It has four edges, and each vertex $a\in N$ has
exactly $d_M(a)$ neighbors in $U$.  By Lemma
\ref{lem:ten-structure}(iv), $\alpha(L)\leq4$, so $\tau(L)\geq6$.  For any
edge $aa'\in E(M)$, the two $U$-neighborhoods must cover $L$ to prevent an
independent four-set in $Q$, but their total size is at most
$d_M(a)+d_M(a')\leq5$.  This is impossible.

It remains that $e(L)=20$.  Then $e(R)=25$, and
Lemma~\ref{lem:ten-structure}(iii) says that $L=T$, the balanced two-fold
blow-up of $C_5$.  Here $A=15$ and $Y\in\{5,6\}$.

If $Y=5$, then $M=\comp{Q[N]}$ has five edges, $X=10$, and the cross-degree
$x_a$ of each $a\in N$ equals $d_M(a)$.  Some edge $aa'$ of $M$ satisfies
\begin{equation}\label{eq:small-degree-sum}
  d_M(a)+d_M(a')\leq5.
\end{equation}
Indeed, otherwise every edge would have endpoint degree sum six and hence
would meet all four other edges.  A pairwise-intersecting family of edges in
a simple graph is contained in a star or a triangle, neither of which has
five edges on five vertices.  The two $U$-neighborhoods of the edge in
\eqref{eq:small-degree-sum} cannot cover $T$, contradicting
$\tau(T)=6$.

If $Y=6$, then $M$ has four edges and $X=9$.  Minimum degree gives
$x_a\geq d_M(a)$ for each $a$, with total excess one; let $z$ carry the
excess.  Every edge of $M$ must meet $z$, since the endpoints of any other
edge have at most five $U$-neighbors and cannot cover $T$.  Hence
$M=K_{1,4}$, centered at $z$, with $x_z=5$ and cross-degree one at each
leaf.  For each center-leaf edge, the two neighborhoods form a minimum
six-vertex cover of $T$ and are disjoint.  These four minimum covers all
contain the fixed five-set $N_Q(z)\cap U$.  Distinct minimum covers of $T$
intersect in at most four vertices, so the covers coincide.  The four leaf
neighborhoods are therefore the same singleton, say $\{u\}$.  The six-set
consisting of $v$, $u$, and the four leaves spans fourteen $Q$-edges: four
from $v$, six among the leaves, and four from $u$.  This violates
admissibility.  We conclude that $\delta(H)\ne4$.

\subsection*{Excluding minimum degree five in the residual graph}

We are left with $\delta(H)=5$.  Choose a degree-five vertex $v$, let
$N=N_H(v)$, let $U$ be its fifteen nonneighbors, put $R=H[U]$, and put
$L=\comp R$.  Then
\begin{equation}\label{eq:H-fifteen}
  e(R)\leq55-5-14=36.
\end{equation}
If $L$ is three-partite, each of its parts has size five.  The graph $R$
then consists of three $K_5$'s with at most six cross-edges.  Every vertex
of $N$ has at most one neighbor in each clique.  If two vertices of $N$
were nonadjacent, each clique would supply three candidates avoiding both.
At most six cross-edges rule out at most eighteen of the twenty-seven
candidate triples, leaving an independent triple.  Together with the two
vertices of $N$, it gives an independent five-set in $H$.  Thus $N=K_5$,
which makes a forbidden $K_6$ with $v$.

Consequently $L$ is not three-partite.  Theorem~\ref{thm:KP} gives
$e(R)\geq34$, and Lemma~\ref{lem:KP-equality}(ii) excludes equality.  Hence
$e(R)$ is $35$ or $36$.

Suppose first that $e(R)=35$.  Then $A=15$ and $Y\in\{5,6\}$.
Lemma~\ref{lem:fifteen-35} gives
\[
  R=C\mathbin{\dot\cup}R_0,
\]
where $C=K_5$ and $T=\comp{R_0}$ is the balanced two-fold blow-up of $C_5$.
Every vertex of $N$ has at most one neighbor in $C$.

If $Y=5$, then $M=\comp{H[N]}$ has five edges, $X=10$, and every
cross-degree equals the corresponding degree in $M$.  Choose an edge
$aa'\in E(M)$ satisfying \eqref{eq:small-degree-sum}, and choose
$c\in C$ avoiding both endpoints.  To prevent $a,a',c$ together with an
uncovered edge of $T$ from forming an independent five-set, the two
$R_0$-neighborhoods must cover $T$.  Their total size is at most five,
contrary to $\tau(T)=6$.

If $Y=6$, then $M$ has four edges and the cross-degrees have one unit of
excess over the degrees in $M$.  The same cover argument forces $M$ to be
$K_{1,4}$ and puts the excess at its center.  For each center-leaf edge, all
six cross-neighbors must lie in $R_0$, since they are needed to cover $T$.
The minimum-cover intersection property again makes the four leaf
neighborhoods a common singleton.  The vertex $v$, that singleton, and the
four leaves span fourteen edges, a contradiction.

Finally suppose that $e(R)=36$.  All counts are equalities:
\[
  A=14,\qquad Y=6,\qquad X=8.
\]
The graph $M=\comp{H[N]}$ has four edges, and each cross-degree is the
corresponding degree in $M$.  Choose any edge $aa'\in E(M)$; its endpoint
degree sum is at most five.

Apply Lemma~\ref{lem:fifteen-36}.  In alternative (i), let $C$ be the
$K_5$, let $c_0$ be the endpoint in $C$ of its unique edge to $R_0$, and
let $T=\comp{R_0}$.  Each of $a,a'$ has at most one neighbor in $C$, so
there is a vertex
\[
  c\in C\setminus\bigl(\{c_0\}\cup N_H(a)\cup N_H(a')\bigr).
\]
The vertex $c$ has no $R$-neighbor in $R_0$.  The two
$R_0$-neighborhoods of $a,a'$ must therefore cover $T$, but their total
size is at most five and $\tau(T)=6$.

In alternative (ii), choose $c$ in the isolated $K_5$ avoiding $a,a'$.
Now $T=\comp{R_0}$ has independence number at most four, so
$\tau(T)\geq6$.  The same uncovered-edge argument again requires a vertex
cover of size at most five.  This final contradiction proves that the graph
$H$ in \eqref{eq:K5H} cannot exist, and completes the proof of
Theorem~\ref{thm:main}.

\section{Concluding remarks}

The proof is confined to five colors.  It does not establish the general
Erd\H{o}s--Gy\'arf\'as conjecture, and the higher-color cases are not addressed
here.  The argument above is entirely graph-theoretic.  In particular, no
finite-search certificate or assumption about a structured coloring family
enters the proof.

\appendix
\section{Dependency audit}

The external inputs are Brooks's theorem and the Kang--Pikhurko extremal
theorem, including its equality description for $(r,n)=(2,10)$ and
$(3,15)$.  All other finite structure statements used above are proved in
the text.

Minimum degree is inherited only after equality in
Lemma~\ref{lem:accounting} gives $X=0$ and peels off an isolated clique.
Every later minimum-degree argument chooses a new minimum-degree vertex in
the graph under consideration.  At order fifteen, the possibility that the
complement is three-partite is treated separately; at order ten, the
possibility that the triangle-free complement is bipartite is also treated
separately.

The main arithmetic identities are
\[
  5\cdot65=325=\binom{26}{2},
  \qquad A\geq14\text{ when }d=5,
\]
and the two near-extremal splits are
\[
  35-4=31,
  \qquad 36-4=32.
\]

For reference, the proof dependencies are as follows.
\begin{enumerate}[label=\textup{D\arabic*.},leftmargin=2.8em]
\item Corollary~\ref{cor:set-color} uses Theorem~\ref{thm:main}, the
color-complement equivalence, and the affine-plane construction on
$\Ffive^2$ for the lower bound.
\item Proposition~\ref{prop:initial} uses the local bounds
\eqref{eq:local}, Brooks's theorem, and Theorem~\ref{thm:KP} with four
parts.
\item The terminal results in Section~\ref{sec:terminal} use only the
local six-set bound, elementary triangle-free arguments, and the two-part
case of Theorem~\ref{thm:KP} with Lemma~\ref{lem:KP-equality}(i).
\item Lemmas~\ref{lem:alpha3-obstruction} and
\ref{lem:alpha3-large} use the preceding terminal lemmas,
Lemma~\ref{lem:accounting}, and the three-part case of
Theorem~\ref{thm:KP}.
\item Lemmas~\ref{lem:fifteen-35} and \ref{lem:fifteen-36} use
Lemma~\ref{lem:ten-structure} and recheck the bipartite complement branch
before applying its nonbipartite alternatives.
\item Lemma~\ref{lem:alpha4-obstruction} uses the independence-three
obstructions and, at its terminal fifteen-vertex stage,
Lemma~\ref{lem:fifteen-35}.
\item Proposition~\ref{prop:equalization} combines
Proposition~\ref{prop:initial}, Lemma~\ref{lem:accounting}, and
Lemma~\ref{lem:alpha4-obstruction}; the identity
$5\cdot65=\binom{26}{2}$ then equalizes all color classes.
\item The final reduction uses the equality cases in the preceding
lemmas, Lemmas~\ref{lem:fifteen-35} and \ref{lem:fifteen-36}, and the
minimum-cover description in Lemma~\ref{lem:ten-structure}.
\end{enumerate}

\bibliographystyle{amsplain}
\bibliography{references}

\end{document}
